\scanpage{162}\mspage{\unsure{17}}
stant -- $\pi_1^{ar}(J)\simeq H^0(k,J)\bigr]$ sous une forme cohomologique, en énonçant que l'on a
\[
\Delta_0\,\mathbb{R}\Gamma_k(J^*)\simeq\mathbb{R}\Gamma_{k}(\Delta J^*[1])
\]
où $\Delta$ est la dualité de Serre, $\Delta_0$ la dualité de \unsure{Pontrjagin|Pontryagin} des groupes abéliens top. totalement disconnexes (dualité à valeurs dans $\mathbb{Q}/\mathbb{Z}$) $J^*$ un complexe de ind-pro-groupes algébriques sur $k$. Tenant compte de cette ``formule de dualité de Lang'' et appliquant $\mathbb{R}\Gamma_{k}$ à la formule de dualité \inserted{des corps de classes} géométrique, on trouve la ``formule de dualité du corps de classes arithmétique''
\[
\boxed{\Delta_0\bigl(H_!(X,F)\simeq H^*\bigl(X,D(F)[3]\bigr)}
\]
(isom de groupes topologiques tot. discontinus).

Autre remarque : lorsque $F$ n'est pas un ``faisceau étale'', mais a une structure continue comme $\alpha_p$, il faut faire attention dans la définition de $Rf_!(F)$, pour $X$ non \unsure{complète|complèt|complet}, à partir de la compactifiée \inserted{$\widetilde X$} ; ainsi, si $F$ provient d'un faisceau \inserted{``admissible''} $\widetilde F$ sur $\widetilde X$, il faut avoir un triangle exact
\diagram{\begin{tikzcd}[ampersand replacement=\&,column sep=small,row sep=small]
\& Rf_!(\widehat{\widetilde F}) \arrow[dl] \& \\
Rf_!(F) \arrow[rr,dash,"{\unsure{\mid|>}}"{pos=1,anchor=center,inner sep=0pt}] \& \& R\widetilde f_{\unsure{*|+}}(\widetilde F)\,, \arrow[ul]
\end{tikzcd}}
% archive sheet number: 162
